\documentclass[11pt,letterpaper]{article}
\usepackage{amsmath}
\title{Notes on Keeping a Small Vegetable Garden}
\author{Marguerite Odell}
\date{March 3, 2024}
\begin{document}
\maketitle

\begin{abstract}
This short note collects practical observations about planning, planting and
maintaining a vegetable garden of modest size. It is written for beginners who
have a sunny corner, a few tools and a considerable amount of enthusiasm.
\end{abstract}

\section{Planning the Beds}
A garden begins on paper. Before any soil is turned, it is worth walking the
site at different hours of the day and noting where the shadows fall. Most
vegetables need at least six hours of direct sunlight, and
\emph{tomatoes}, peppers and squash greatly prefer more.

\subsection{Choosing a Location}
Level ground is convenient, but it is not essential. Water drains poorly in
hollows, so an apparently ideal flat patch may turn into a swamp after the
first heavy rain.\footnote{A simple test is to dig a hole, fill it with water,
and measure how long it takes to empty.}

\subsection*{A Word on Soil}
Good soil is dark, crumbly and smells faintly sweet. Compost improves almost
every kind of soil, whether sandy or heavy with clay.

\section{What to Plant}
The following crops are forgiving and productive:
\begin{itemize}
\item Radishes, which are ready in about a month.
\item Lettuce and other leafy greens, including
  \begin{itemize}
  \item spinach in the cool seasons,
  \item chard, which survives almost any neglect.
  \end{itemize}
\item Beans, both bush and climbing varieties.
\end{itemize}

A rough weekly routine might run as follows.
\begin{enumerate}
\item Inspect the plants for pests and disease.
\item Water deeply, preferably in the early morning.
\item Remove weeds before they set seed:
  \begin{enumerate}
  \item pull small ones by hand,
  \item hoe the larger ones between rows.
  \end{enumerate}
\item Harvest whatever is ripe.
\end{enumerate}

\section{Spacing}
Seed packets give spacing in inches, but a useful rule is that a bed of width
$w$ feet and length $\ell$ feet holds roughly $n \approx 4w\ell/s^2$ plants when
the spacing is $s$ inches, scaled by $144/4$. Overcrowded plants compete for
light, water and nutrients, and disease spreads rapidly among them.

\bigskip
\begin{center}
Happy gardening, and may the frost arrive late.
\end{center}
\end{document}
